% RESEARCH_DRAFT_SHA256: 6499d458fbf38b86760184cd4456df82dc5af0c9d649d83127aa643465648ee7
% Layout and source-reference wrappers expanded for the web renderer.
\documentclass[11pt]{article}
\usepackage[margin=25mm]{geometry}
\usepackage{fontspec,unicode-math}
\setmainfont{Latin Modern Roman}
\setmathfont{Latin Modern Math}
\usepackage{amsmath,amssymb,mathtools,xurl,hyperref}
\hypersetup{colorlinks=true,urlcolor=blue,linkcolor=blue}
\setcounter{secnumdepth}{0}
\setlength{\parindent}{0pt}
\setlength{\parskip}{5pt}
\setlength{\emergencystretch}{3em}
\newcommand{\sref}[2]{\hyperlink{source:#1:#2}{[S#2]}}
\newcommand{\cataloguescope}[1]{\paragraph{Source record.}#1}
\newcommand{\floor}[1]{\left\lfloor#1\right\rfloor}
\begin{document}
\section*{Research-attempt provenance}
This standalone research attempt retains the complete original definition body below,
including its historical source and preparation notes. Those notes describe the
earlier catalogue preparation; the new attempt and its verification are identified
separately after that body. No original ID, catalogue status, or source statement is changed.

% UnsolvedMath ID 30000093; source catalogue code OWR-735-002
\section{Fourier Reconstruction of Discrete Objects}\label{30000093}
{\small UnsolvedMath ID 30000093 · Computer Science}
\subsection{Definitions and mathematical statement}

\paragraph{Shared notation.}Standard number systems are $\mathbb N,\mathbb Z,\mathbb Q,\mathbb R,\mathbb C$. Specialized notation and index ranges follow the cited source and the definitions below.


For a function on a finite grid, discrete Fourier coefficients are linear combinations of its values with complex exponential weights. Observing only frequencies in $\Omega$ generally leaves multiple possible functions unless further structure is imposed.
\paragraph{Record scope and qualifications.}
The catalogue title is retained as metadata. Named certificates, contracts, bounds or frameworks in that title are not asserted as results of the original source extract.
The supplied registry flags this extraction as unrecoverable. It is displayed as a literal source fragment; the definitions provide context rather than a replacement formal statement.

The registry flags this extraction as unrecoverable, and the supplied literature assessment refers to a different context. No unique intended reconstruction theorem is inferred.

Source-record limitation: Source flags the statement as unrecoverable: The cited official report on fast methods for non-local operators contains no discrete Fourier-reconstruction problem; the statement belongs to the following Oberwolfach report. No clean statement is promoted under the attached source, and the original remains retained for provenance.
The source text below is rendered literally where necessary; damaged formulas and composite extracts are not silently repaired into a new mathematical claim.
The supplied catalogue marks this record as unrecoverable. Its extract and terminology are documented, but a complete original problem has not been reconstructed.
\paragraph{Statement recorded in the supplied source.}
Question: Under what assumptions can a discrete object f be reconstructed from its Fourier coefficients \textbackslash{}hat f(k) for k\textbackslash{}in\textbackslash{}Omega\textbackslash{}subset\textbackslash{}\{0,\textbackslash{}ldots,N-1\textbackslash{}\}\textasciicircum{}d? \href{https://doi.org/10.4171/owr/2004/33}{[S2]}
\paragraph{Source status.}
The supplied catalogue labels this source record partially solved; that metadata was checked against the supplied snapshot on 2026-10-11. The label is preserved for provenance and is not an independent verification of present open status or mathematical correctness. The catalogue literature assessment is dated 2026-08-21.
\subsection{Short English statement}
Ask which assumptions permit reconstruction from selected Fourier coefficients.
\subsection{Sources}
\begin{itemize}
\item[S1] Ulam AI and the UnsolvedMath Contributors, UnsolvedMath: A Curated Collection of Open Mathematics Problems, record 30000093 (OWR-735-002). Catalogue curation and metadata: CC BY 4.0. The source extract is retained; definitions, English summaries and qualifications are editorial additions.
\url{https://huggingface.co/datasets/ulamai/UnsolvedMath}.
\item[S2] Fast Numerical Methods for Non-local Operators (2005), Oberwolfach Reports, DOI 10.4171/owr/2004/33, pp. 1747--1788
\url{https://doi.org/10.4171/owr/2004/33}.
\item[S3] Report full text
\url{https://ems.press/content/serial-article-files/45951}.
\item Source-specific Oberwolfach Reports licences, including CC BY-SA 4.0 where applicable, are retained; see THIRD\_PARTY\_NOTICES.md.
\end{itemize}
\paragraph{Preparation and review.}
The original source and UnsolvedMath contributors retain attribution for the source extract. Codex (GPT-6; exact serving revision not exposed) assisted with definitions, English summaries, scope notes and formatting. No new human mathematical authorship or independently verified solution is asserted. Permanent source ID 30000093 and code OWR-735-002 checked against the complete supplied registry and current repository. Distinct upstream records, including repeated questions and broad or damaged extracts, are retained for complete coverage of the source CS category. A repeated record is not a new independent conjecture. Complete mathematical-equivalence and current-literature review has not been performed. Preparation and final review follow contribution prompts C and D, with source gaps disclosed. Structural and rendering checks do not establish a new theorem or open status.
\label{end:30000093}

\section{Literature and context check}
The cached primary report 33/2004 was inspected: its title and introduction concern Fast Numerical Methods for Non-local Operators. The supplied unrecoverable/misattributed-source flag is retained; this inspection does not recover a unique original Fourier-reconstruction theorem. The explicit one-dimensional sparse model below is a classical Prony/Vandermonde route, rederived to test what a structural assumption can accomplish. It is not silently adopted as the source's unknown complete target.
\section{Research attempt: Exact sparse reconstruction and the missing source target}
On $\mathbb Z/N\mathbb Z$, set $\zeta=e^{-2\pi i/N}$ and define the unnormalized Fourier data $m_k=\sum_{x=0}^{N-1}f(x)\zeta^{kx}$. Without a structural restriction, missing a frequency $k_0$ precludes uniqueness: the nonzero vector $f(x)=e^{2\pi i k_0x/N}$ has $m_k=0$ for every $k\ne k_0$. This follows by summing the finite geometric series whose ratio is an $N$th root of unity different from one. Hence it is indistinguishable from zero on any observed set omitting $k_0$.

Now assume exactly $s\geq1$ nonzero entries and observe $m_0,\ldots,m_{2s-1}$, with $2s\leq N$. Write their distinct locations as $x_1,\ldots,x_s$, their nonzero values as $c_1,\ldots,c_s$, and $z_j=\zeta^{x_j}$. The $s$ by $s$ Hankel matrix $H=(m_{a+b})_{a,b=0}^{s-1}$ factors as
\[
 H=V\operatorname{diag}(c_1,\ldots,c_s)V^T,
 \qquad V_{a,j}=z_j^a.
\]
The Vandermonde determinant $\det V=\prod_{i<j}(z_j-z_i)$ is nonzero: the determinant is alternating, vanishes at each equality of two variables, has the product's degree and leading coefficient, proving the formula. Thus $H$ is invertible even though it need not be Hermitian or positive definite.

Solve the linear system
\[
 \sum_{b=0}^{s-1}m_{a+b}q_b=-m_{a+s},\qquad 0\leq a<s.
\]
The polynomial $Q(z)=z^s+\sum_{b=0}^{s-1}q_bz^b$ is exactly $\prod_j(z-z_j)$. Indeed the latter polynomial satisfies the displayed system after substituting $m_k=\sum_jc_jz_j^k$, and invertibility proves uniqueness. Evaluate $Q$ at all $N$ roots $\zeta^x$ to recover the support, and then solve the Vandermonde system for $c_j$ from the first $s$ moments. This is an explicit exact-arithmetic reconstruction using $O(s^3+Ns)$ scalar arithmetic operations, with conventional elimination and direct evaluation. It is not a finite-precision stability theorem or a bit-complexity bound for arbitrary complex inputs.

Uniqueness under an at-most-$s$ promise can also be proved without first knowing the exact support size. The difference of two such vectors has at most $2s$ nonzero entries. If their observed moments agree, the first $t$ equations, where $t\leq2s$ is the difference's actual support size, form a square Vandermonde system and force all entries to vanish.

The consecutive-frequency hypothesis is substantive. With even $N$ and only even frequencies, the point masses at $0$ and $N/2$ have identical data, since $\zeta^{kN/2}=1$ for every even $k$. Even one-sparse objects need not be identifiable from an arbitrary set of the same number of frequencies.
\section{Verification and adversarial audit}
The dense obstruction, exactly sparse constructive proof, and at-most-sparse uniqueness are distinct statements. The zero object is handled by the at-most-sparse argument. The Hankel proof uses nonzero coefficients and distinct nodes, and its transpose is deliberately not replaced by a conjugate transpose. An explicit even-frequency alias tests the measurement hypothesis.
\paragraph{Audit corrections.} The supplied wrong-source/unrecoverable flag was preserved. A complete proof in an auxiliary sparse model was not labeled a recovered or solved original problem.
\paragraph{External computation.} None; symbolic arguments and boundary checks only.
\section{Scope and remaining gaps}
No claim is made for arbitrary dimensions, arbitrary frequency sets, noisy data, binary-valued objects, or the unrecovered original source formulation. The exact sparse theorem is known material and leaves the original damaged record unresolved.
\paragraph{Integration review.} Codex (GPT-6; exact serving revision not exposed) reviewed the written arguments and expanded layout and source-reference wrappers. Repository metadata and the production web renderer were checked using Python 3.13 and JavaScript with MathJax 3.2.2. This model-assisted review is not independent human proof verification; no PDF compilation or formal proof checking is asserted.

\section{Completion Estimate}
\noindent\textbf{COMPLETION: 0\%}
The calculations below the preserved source statement investigate a concrete route,
but establish no substantiated new progress toward the entire original target.
Known special cases and failed sufficient conditions are not credited as a solution.
\section{References}
The original source references and catalogue attribution are preserved in the
Sources section above. The following URLs were actually consulted for this attempt;
retention of an additional original reference does not claim fresh consultation of it.
\begin{itemize}
\item \url{https://ems.press/content/serial-article-files/45951} (consulted 2026-10-11). 
\end{itemize}
\section{Model and contributor disclosure}
OpenAI gpt-6-astra, requested reasoning configuration ultra, generated this research
attempt on 2026-10-11. The exact serving revision is not exposed. Codex GPT-6
performed orchestration, repository-format review, and written mathematical review at the declared scope; its exact serving revision
is also not exposed. The mathematical route, boundary checks and final adversarial
self-audit followed contribution prompts A then B in
\texttt{docs/contribution\_prompts.md}, after reading \texttt{CONTRIBUTING.md}.
Python 3.14.2 assembled the standalone documents; any mathematical computation is
specified in Verification. No theorem prover was used. No human mathematical
authorship, substantial human proof correction, or independent human proof checking
is asserted. Model-assisted review, source retrieval, and structural checks do not
establish correctness or novelty. Full private reasoning transcripts are unavailable
as reproducibility artifacts; the complete arguments, inputs and outcomes actually
used are stated here or in the companion batch audit.
\section{Final status}
Unresolved as a new research contribution. This is a concrete route investigation,
not a claim of new complete resolution or a current-status certification. A known
answer reproduced for context does not become a new contribution. The catalogue's
source-status label is preserved; this attempt status does not assert that a
source-labelled solved question is open.
\end{document}
